lastly, we have · -- Case C odd: both order p² → C_{p²} ⋊ C_p right sorry

we need to create an OddCaseC.lean file just like we did for the OddCaseB.lean file for example.
I can do the linking to classification after (filling in the sorry).

here is the proof skeleton.

Case C: a and b both have order p2 .

By Fact (9), we can show that for odd prime, it is possible to make a substitution and get into Case B. PLACEHOLDER FOR INFORMATION TO BE FILLED IN: [SHOW MORE]

Fact 9 appears to be this:

Formula for powers of product in group of class two

Suppose G {\displaystyle G} is a group and x , y ∈ G {\displaystyle x,y\in G} are elements such that the subgroup generated by x {\displaystyle x} and y {\displaystyle y} has nilpotence class two. Equivalently, suppose that [ x , y ] {\displaystyle [x,y]} commutes with both x {\displaystyle x} and y {\displaystyle y}, where we define:

[ x , y ] = x − 1 y − 1 x y {\displaystyle [x,y]=x^{-1}y^{-1}xy}

Then, we have the identity:

x r y r = [ x , y ] r ( r − 1 ) / 2 ( x y ) r {\displaystyle x^{r}y^{r}=[x,y]^{r(r-1)/2}(xy)^{r}}

(Note that the formula remains the same even if we follow the other definition of commutators, i.e.:

[ x , y ] = x y x − 1 y − 1 {\displaystyle [x,y]=xyx^{-1}y^{-1}}

because, in this case, the two commutators are the same.) Proof

To go back and forth between x r y r {\displaystyle x^{r}y^{r}} and ( x y ) r {\displaystyle (xy)^{r}}, we need to flip adjacent x {\displaystyle x}'s and y {\displaystyle y}'s. Each flip involves a multiplication by [ x , y ] {\displaystyle [x,y]}. Explicitly, we can rewrite:

y x = [ x , y ] − 1 x y {\displaystyle yx=[x,y]^{-1}xy}

The expression [ x , y ] − 1 {\displaystyle [x,y]^{-1}} is central because the group has class two, so it can be pulled to the left. The number of such flip operations we need to do is 0 + 1 + 2 + ⋯ + ( r − 1 ) {\displaystyle 0+1+2+\dots +(r-1)} (0 for the right-most y {\displaystyle y}, 1 for the second y {\displaystyle y} from the right, and so on), which becomes r ( r − 1 ) / 2 {\displaystyle r(r-1)/2}. Thus, we get:

( x y ) r = ( [ x , y ] − 1 ) r ( r − 1 ) / 2 x r y r {\displaystyle (xy)^{r}=([x,y]^{-1})^{r(r-1)/2}x^{r}y^{r}}

Rearranging gives us the formula in the suggested form.

The detailed manipulations are shown for r = 2 {\displaystyle r=2} and r = 3 {\displaystyle r=3} below. Case r = 2 {\displaystyle r=2}

In this case:

( x y ) 2 = x y x y = x ( y x ) y = x ( [ x , y ] − 1 x y ) y = [ x , y ] − 1 x ( x y ) y = [ x , y ] − 1 x 2 y 2 {\displaystyle (xy)^{2}=xyxy=x(yx)y=x([x,y]^{-1}xy)y=[x,y]^{-1}x(xy)y=[x,y]^{-1}x^{2}y^{2}}

Rearranging gives:

x 2 y 2 = [ x , y ] ( x y ) 2 {\displaystyle x^{2}y^{2}=x,y^{2}}

This fits the formula, because for r = 2 {\displaystyle r=2}, r ( r − 1 ) / 2 = 1 {\displaystyle r(r-1)/2=1}. Case r = 3 {\displaystyle r=3}

In this case:

( x y ) 3 = x y x y x y = x y x ( y x ) y = x y x ( [ x , y ] − 1 x y ) y = [ x , y ] − 1 x y x 2 y 2 = [ x , y ] − 1 x ( y x ) x y 2 = [ x , y ] − 1 x ( [ x , y ] − 1 x y ) x y 2 {\displaystyle (xy)^{3}=xyxyxy=xyx(yx)y=xyx([x,y]^{-1}xy)y=[x,y]^{-1}xyx^{2}y^{2}=[x,y]^{-1}x(yx)xy^{2}=[x,y]^{-1}x([x,y]^{-1}xy)xy^{2}}

= [ x , y ] − 2 x 2 y x y 2 = [ x , y ] − 2 x 2 ( y x ) y 2 = [ x , y ] − 2 x 2 ( [ x , y ] − 1 x y ) y 2 = [ x , y ] − 3 x 3 y 3 {\displaystyle =[x,y]^{-2}x^{2}yxy^{2}=[x,y]^{-2}x^{2}(yx)y^{2}=[x,y]^{-2}x^{2}([x,y]^{-1}xy)y^{2}=[x,y]^{-3}x^{3}y^{3}}

Rearranging gives:

x 3 y 3 = [ x , y ] 3 ( x y ) 3 {\displaystyle x^{3}y^{3}=[x,y]^{3}(xy)^{3}}

This fits the formula because for r = 3 {\displaystyle r=3}, r ( r − 1 ) / 2 = 3 ( 3 − 1 ) / 2 = 3 {\displaystyle r(r-1)/2=3(3-1)/2=3}.